\section{The contraction region in the shift--length plane}\label{sec:region}

The route this program must serve is Proposition~\ref{prop:route}: contraction
along a sequence with $L_j\to\infty$ and $\omega_j\downarrow0$. It is natural to
ask what the region
\[
 \mathcal R=\big\{(\omega,L):\ \lVert V_{\omega,L}\rVert\leq1\big\}
\]
looks like, and the earlier shifted-zeta investigation
\cite{CriticalPath,CumulativePath} looked for a \emph{critical path} inside it.
Proposition~\ref{prop:unimodular} settles the question completely, and the answer
is that the region has no shape worth computing --- which is itself informative,
and which converts the plane into an instrument of a different kind.

\subsection{The region is monotone, and its supremum is one or infinity}

\begin{proposition}[Monotonicity]\label{prop:monotone}
For each fixed $\omega$, $L\mapsto\lVert V_{\omega,L}\rVert$ is nondecreasing, and
$\sup_L\lVert V_{\omega,L}\rVert=\lVert V_\omega\rVert_{L^2(\R)}$.
\end{proposition}

\begin{proof}
$V_{\omega,L}=P_LV_\omega P_L$ and the intervals are nested, so for $L<L'$ the
first is a compression of the second. $P_L\to I$ strongly, so
$\lVert P_LV_\omega P_Lf\rVert\to\lVert V_\omega f\rVert$ for each $f$.
\end{proof}

\begin{proposition}[Dichotomy]\label{prop:dichotomy}
Fix $\omega\in(0,\frac12]$ and suppose no zero of $\xi$ has
$|\re\rho-\frac12|=\omega$ exactly. Then
\[
 \sup_L\lVert V_{\omega,L}\rVert=
 \begin{cases}
 1,&\text{if $\xi$ has no zero with }|\re\rho-\tfrac12|>\omega,\\[2pt]
 \infty,&\text{otherwise.}
 \end{cases}
\]
\end{proposition}

\begin{proof}
The poles of $K_\omega$ are at $p=\rho-\frac12-\omega$, so one lies in
$\re p>0$ exactly when some $\re\rho>\frac12+\omega$; by the functional equation
that is the same as some $|\re\rho-\frac12|>\omega$. If there is none, the
argument of Proposition~\ref{prop:flux} applies verbatim with RH replaced by that
weaker hypothesis: $K_\omega$ is inner and $V_\omega$ is unitary, of norm exactly
one. If there is one, at $\re p=\delta-\omega>0$, moving the contour picks up a
residue and $k_\omega$ acquires a component growing like $e^{(\delta-\omega)x}$,
so $V_\omega$ is unbounded on $L^2(\R)$; by Proposition~\ref{prop:monotone} the
truncations are then unbounded in $L$.
\end{proof}

\begin{corollary}[The shifted family is a graded criterion]\label{cor:graded}
For each $\omega\in(0,\frac12]$ the following are equivalent:
\begin{enumerate}
\item $\lVert V_{\omega,L}\rVert\leq1$ for every $L>0$;
\item $\xi$ has no zero with $|\re\rho-\frac12|>\omega$.
\end{enumerate}
In particular, under the Riemann hypothesis $\mathcal R$ is the whole quadrant
$(0,\frac12]\times(0,\infty)$.
\end{corollary}

So the family is not one criterion but a totally ordered family of them, one per
shift, interpolating between the trivial statement at $\omega=\frac12$ --- no
zeros outside the critical strip --- and the Riemann hypothesis as
$\omega\downarrow0$. \textbf{Each $\omega$ buys exactly a zero-free strip of width
$\omega$.}

\begin{remark}[The two limits of Proposition~\ref{prop:route} do different jobs]
$L_j\to\infty$ realizes the supremum of Proposition~\ref{prop:monotone} and
covers every test function; $\omega_j\downarrow0$ closes the detection threshold
of Corollary~\ref{cor:graded}. Neither substitutes for the other.
\end{remark}

\subsection{There is no critical path, and what the plane is for}

Corollary~\ref{cor:graded} says $\mathcal R$ has no interesting boundary. Under
the Riemann hypothesis it is everything; if the hypothesis fails at distance
$\delta$, it is everything for $\omega>\delta$ and cut off at a finite $L$ for
$\omega<\delta$. The $L$-direction is monotone bookkeeping and the entire content
sits at $L=\infty$, where it is the $H^\infty$ norm of $K_\omega$. There is no
curve in the plane whose shape is worth computing.

This also answers the question the manuscript previously left open, whether the
flow formulation escapes the saturation constraint of \cite{Selection}. It does
not, and the reason is at the level of quantifiers: a mechanism producing a
strict contraction at one fixed $\omega>0$, however uniform in $L$, would
establish a zero-free strip of width $\omega$ --- a real and currently open
theorem, and a strong one --- but not the Riemann hypothesis. The quantifier
missing from such a statement is $\omega\downarrow0$, and by
\eqref{eq:absorption} the margin vanishes with $\omega$ by construction.

Monotonicity does buy one thing, and it is worth stating on its own.

\begin{remark}[A disproof instrument]\label{rem:disproof}
A finite computation exhibiting $\lVert V_{\omega,L}\rVert>1$ at a single
$(\omega,L)$ proves that $\xi$ has a zero at distance more than $\omega$ from the
critical line. A computation showing $\lVert V_{\omega,L}\rVert\leq1$ at any
finite $L$ proves nothing whatever. Numerics on the shift--length plane are
therefore one-sided in the same way the companion investigation's disproof test
is, and should be organized as such: certified \emph{lower} bounds on Rayleigh
quotients, scanning $\omega$ downward and $L$ upward.
\end{remark}

\subsection{The generator region, and what controls its boundary}

The generator condition $Q_{\omega,L}\succeq0$ is strictly stronger than
contraction --- \cite{CriticalPath} exhibits a certified input at
$L=\log7$, $\omega=10^{-11}$ on which the generator is negative while the
cumulative defect is positive --- and it \emph{does} have a boundary. That
boundary is controlled by
$\kappa_L=\lVert X\rVert_{Q\to Q}$, through the sufficient path condition
$\omega\kappa_L<\pi/4$ of \cite{CriticalPath}. Proposition~\ref{prop:hyperbolic}
and the spectral form identify what $\kappa_L$ measures.

\begin{proposition}[$\kappa_L$ is an inverse zero-placement precision]
\label{prop:kappa}
With $X_af=(x-a)f$ one has $\widehat{X_af}=i\widehat F'-a\widehat F$, so by
\eqref{eq:spectral-target}
\[
 \kappa_L^2=\sup_f\frac{\sum_\rho\big|i\widehat F'(\gamma_\rho)
 -a\widehat F(\gamma_\rho)\big|^2}{\sum_\rho|\widehat F(\gamma_\rho)|^2}
 \;\xrightarrow[\text{near-null }f]{}\;
 \sup_f\frac{\sum_\rho|\widehat F'(\gamma_\rho)|^2}
 {\sum_\rho|\widehat F(\gamma_\rho)|^2},
\]
the ratio of derivative samples to value samples of $\widehat F$ at the Riemann
zeros. If $e_L$ is the ground vector of $Q_{0,L}$ and $\varepsilon_L$ the
root-mean-square distance from the zeros of $\widehat e_L$ to the nearest
$\gamma_\rho$, then $m_L\approx\varepsilon_L^2Q[Xe_L]$ and
$\kappa_L\approx\varepsilon_L^{-1}$, so
$\kappa_L\sqrt{m_L}\approx\sqrt{Q[Xe_L]}$.
\end{proposition}

The identity for $\kappa_L^2$ is exact; the relation
$\kappa_L\sqrt{m_L}\approx\sqrt{Q[Xe_L]}$ is a single-parameter heuristic. It is
tested by the values recorded in \cite{CriticalPath} across twenty orders of
magnitude in $m_L$:

\begin{center}
\begin{tabular}{@{}lrrrrr@{}}
\toprule
horizon & $m_L$ & $\kappa_L$ & $\kappa_L\sqrt{m_L}$ & $\pi/(4\kappa_L)$
& recorded root\\
\midrule
$\log3$ & $5.537\times10^{-8}$  & $6.112\times10^{1}$  & $0.01438$
 & $1.29\times10^{-2}$ & $1.64\times10^{-2}$\\
$\log5$ & $9.293\times10^{-18}$ & $2.534\times10^{6}$  & $0.00773$
 & $3.10\times10^{-7}$ & $3.95\times10^{-7}$\\
$\log7$ & $6.802\times10^{-28}$ & $2.326\times10^{11}$ & $0.00607$
 & $3.38\times10^{-12}$ & $4.30\times10^{-12}$\\
\bottomrule
\end{tabular}
\end{center}

$\kappa_L\sqrt{m_L}$ moves by a factor of $2.4$ while $m_L$ moves by $10^{20}$,
and the recorded diagnostic root sits within a factor of $1.3$ of
$\pi/(4\kappa_L)$ throughout. So the generator boundary is
$\omega_*(L)\asymp\sqrt{m_L}$.

\begin{remark}[The generator boundary is an artifact]\label{rem:artifact}
$m_L=\lambda_{\min}(Q_{0,L})$ is positive precisely when Weil positivity holds on
$I_L$, so its superexponential collapse happens \emph{in the world the program is
trying to prove}. The generator region's boundary therefore collapses under the
Riemann hypothesis as well, and carries no arithmetic information; it is a
property of a sufficient condition.

How far is it from anything useful? By Proposition~\ref{prop:dichotomy}, a zero at
distance $\delta$ becomes visible on $I_L$ only once $(\delta-\omega)L\gtrsim1$,
so the working scale is $\omega\sim1/L$, a power law. Against
$\pi/(4\kappa_L)$ the ratios are $7.1\times10^{1}$, $2.0\times10^{6}$ and
$1.5\times10^{11}$ at the three horizons above --- eleven orders of magnitude at
$L=\log7$, widening superexponentially. The sufficient condition and the needed
one are not close and are diverging, which is the quantitative reason the earlier
investigation's elementary extension lemma asked for extension lengths of order
$e^{-\mathrm{const}/m_L}$.
\end{remark}

\subsection{The Cayley coordinate is the positive-real criterion}

\cite{CumulativePath} introduced $Z_{\omega,L}=(I-V)(I+V)^{-1}$ and
$P=\frac2\omega\re Z$, with $D_{\omega,L}=\frac\omega2(I+V^*)P(I+V)$, so that
$P\succeq0\iff\lVert V\rVert\leq1$, and the corresponding causal symbol
$a_{\rm cum,\omega}=\frac2\omega(1-K_\omega)/(1+K_\omega)$.
Proposition~\ref{prop:unimodular} identifies it.

\begin{proposition}[Reactance]\label{prop:cayley}
On the critical line,
\[
 \frac{1-K_\omega(i\tau)}{1+K_\omega(i\tau)}=i\tan\Psi(\tau,\omega),
\]
purely imaginary. Consequently $\lVert V_{\omega,L}\rVert\leq1$ for all $L$ is
equivalent to $Z_\omega$ being accretive, that is to $a_{\rm cum,\omega}$ being a
\emph{positive-real} function on the right half-plane, which by
Corollary~\ref{cor:graded} is equivalent to $K_\omega$ being inner.
\end{proposition}

\begin{proof}
Substitute $K_\omega(i\tau)=e^{-2i\Psi}$ from \eqref{eq:phase}; the Cayley
transform carries the unit circle to the imaginary axis. The equivalences are
Corollary~\ref{cor:graded} together with the congruence of
\cite{CumulativePath}.
\end{proof}

That is the classical passivity criterion, and positive-real functions are
Nevanlinna functions after rotation, with a Herglotz representation whose measure
is supported on the resonances --- which is Proposition~\ref{prop:resonance} seen
from the other side. Two cautions. First, Remark~\ref{rem:trap} applies again:
$a_{\rm cum,\omega}$ is purely imaginary on the imaginary axis, so any expansion
carried out there loses the entire form, and the $\omega^2$ expansion of
\cite{CumulativePath} is safe only because it is taken on a far-right Laplace
line. Second, the natural tool for such a function is its Herglotz
representation, not a Taylor expansion in the shift. Both point at the same
literature as the companion investigation's top open item \cite{SuzukiDeBranges}.

\subsection{What the margin is}

Combining \eqref{eq:absorption} with the spectral form gives the margin an exact
reading:
\[
 \lambda_{\min}(Q_{0,L})
 =\min_{\lVert f\rVert=1,\ \supp f\subseteq I_L}\ \sum_\rho|\widehat F(\gamma_\rho)|^2 ,
\]
that is, \textbf{how well a function supported in an interval can avoid the
Riemann zeros in frequency.} Two consequences.

The band edge is explained rather than observed. The cheapest way to avoid every
$\gamma_\rho$ is to live below the lowest one, so near-null directions are
bandlimited below $\gamma_1=14.1347$ --- which \cite{Bottleneck} recorded as a
measurement and \cite{Selection} carries as the second half of its seventh
selection rule. Here it is a consequence.

But bandlimiting is only the first factor. The recorded margins give
$d\log m_L/dL\approx-44$ and $-69$ on the two available intervals, far steeper
than the bandwidth-$\gamma_1$ prolate rate $-\gamma_1$ \cite{SlepianPollak}. By
Proposition~\ref{prop:kappa} the second factor is the zero-placement precision:
$m_L\approx\varepsilon_L^2Q[Xe_L]$. \textbf{How $\varepsilon_L$ collapses is the
mechanism question}, and it is the form in which this paper leaves the companion
investigation's request for a mechanism rather than a measurement.
